\section{Proposed simulation algorithm for lattice QCD}

With respect to the QCD simulation algorithms used to date,
the combination of the transformations obtained
through the Euler integration of the Wilson flow
and the HMC algorithm is expected to
sample the topological sectors more quickly
and to be generally more efficient.
The use of the Wilson flow is suggested
if the gauge action coincides with the Wilson
plaquette action. For other actions,
the appropriate flow can be easily constructed
following the lines of sect.~4.


\subsection{Choice of the parameters}

The parameters of the Euler integrator
are the integration step size $\eps$ and the number $n$ of
Euler sweeps that are applied. One also needs to choose
a definite ordering of the links of the lattice.

The step size $\eps$ should be positive
and not larger than, say, $1/16$ so that
the in\-ver\-ti\-bi\-lity of the
Euler integrator is guaranteed within a safe margin.
Some tuning of the integration time $n\eps$
will certainly be required in order to maximize
the efficiency of the algorithm.
Note that the unit of time differs
from the one used in subsect.~4.4, i.e.\ setting $t=1$ there
corresponds to
an integration time $n\eps=\beta/32$.

The ordering of the links can be chosen arbitrarily.
One may, for example, first visit the links $(x,0)$
on all even points $x$, then the links $(x,0)$ on
all odd points, then the links $(x,1)$ on all even points,
and so on.
This ordering is well suited for parallel processing,
since the link variables in a given direction on the even (odd) sites
are decoupled from one another and can therefore be updated in parallel.


\subsection{Force calculation}

At the beginning of an update cycle,
the current gauge field
$U$ is transformed to the field $V=\transne^{-1}(U)$ by applying
$n$ backward Euler sweeps to $U$.
The force that drives the molecular-dynamics evolution of $V$
is then given by
\begin{equation}
  F(z,\rho)^c=\hat{\du}^c_{z,\rho}\left\{
  S(\transne(V))-\ln\det\transnes(V)\right\},
\end{equation}
where the action $S(U)$ now includes the usual sea-quark pseudo-fermion actions
(for simplicity, the dependence on the
pseudo-fermion fields is suppressed).

\begin{figure}[htbp]
  \centering
  \includegraphics[width=3.6cm]{images/trj.pdf}
  \caption{%
The proposed algorithm evolves the field $V$ using the standard
HMC algorithm (thick line). The force that drives the
molecular-dynamics evolution
is obtained by forward integration of the Wilson flow
and subsequent backward propagation of
the derivatives of the action and the Jacobian
of the flow ($n=3$ in this figure).
}
  \label{fig:trj}
\end{figure}

Each time the force is to be calculated,
the current field $V$ must be transformed
to $U$ again by applying $n$
forward Euler sweeps (see fig.~3). The fields
$U_0,U_{\eps},\ldots,U_{n\eps}$ generated in this process should
be stored in memory so that they will be available
when the force is propagated from $U$ to $V$.
Note that the intermediate fields
\begin{equation}
  U_{k\eps,[x,\mu]}(y,\nu)=
  \begin{cases}
    U_{k\epsilon}(y,\nu) & \text{if $(y,\nu)\geq(x,\mu)$,}\\[0.8ex]
    U_{k\epsilon+\epsilon}(y,\nu) & \text{otherwise,}
  \end{cases}
\end{equation}
are then also available.

The factorization (5.14) of the Jacobian implies
\begin{equation}
   F(z,\rho)^c=\hat{\du}^c_{z,\rho}S(U_{n\eps})
   -\sum_{k=0}^{n-1}\sum_{x,\mu}
  \hat{\du}^c_{z,\rho}\ln\det\euler{x,\mu,*}(U_{k\eps,[x,\mu]}).
\end{equation}
Moreover,
\begin{equation}
  \hat{\du}^c_{z,\rho}S(U_{n\eps})=
  \sum_{y,\nu}\left.\du^a_{y,\nu}S(U)\right|_{U=U_{n\eps}}
  \transnes(y,\nu;z,\rho)^{ac},
\end{equation}
and there is a similar formula for the other terms in eq.~(6.3)
involving the Jacobian matrices of the
transformations $V\to U_{k\eps,[x,\mu]}$.
All these matrices, as well as the one in eq.~(6.4),
are products of the Jacobian matrices of the one-link
transformations (5.8).
The force can therefore be computed recursively
as follows:

\begin{enumerate}[leftmargin=2.2em,itemsep=0.6ex]
\item Set $U=U_{n\eps}$ and $F(z,\rho)^c=\du^c_{z,\rho}S(U)$.

\item For $k$ from $n-1$ to $0$,
run through all links $(x,\mu)$ in reverse order,
set $U=U_{k\eps,[x,\mu]}$
and update the force according to
\begin{equation}
   F(z,\rho)^c\to\sum_{y,\nu}F(y,\nu)^a
   [\euler{x,\mu,*}(U)](y,\nu;z,\rho)^{ac}
   -\du^c_{z,\rho}\ln\det\euler{x,\mu,*}(U).
\end{equation}
\end{enumerate}

\noindent
Note that the field $U$ backtracks
the forward integration of the Wilson flow
in the course of the recursion.
Since the Jacobian matrix $\euler{x,\mu,*}(U)$ differs from unity
only on the links sharing a plaquette with $(x,\mu)$,
the total computational effort required
for the propagation (6.5) of the force
is expected to be similar to the one required for
$n$ applications of a nearest-neighbour
gauge-covariant difference operator to the force field.


\subsection{Domain-decomposed algorithm}

Field transformations
can also be combined with the DD-HMC algorithm~\cite{DDHMC}.
Only the so-called active link variables
are transformed in this case, but the transformations may depend on
the inactive field components.

In the pure gauge theory, it is then again possible to
construct trivializing flows. They operate
on the active link variables and
contract the gauge action in each do\-main to a constant
(i.e.\ an expression depending on the inactive link variables
only).
These flows are not the same as the ones constructed
in sect.~4, but can be obtained by solving
the differential equations derived there.
In particular,
the trivializing flow
for the Wilson action is, to leading order,
a slightly modified Wilson flow,
where the plaquettes containing $\nu$ active links
are given the weight $4/\nu$.
